\documentclass[11pt,oneside]{report}

\usepackage[margin=0.95in,headheight=24pt]{geometry}
\usepackage[T1]{fontenc}
\usepackage[utf8]{inputenc}
\usepackage{amsmath,amssymb}
\usepackage{booktabs}
\usepackage{array}
\usepackage{xcolor}
\usepackage{hyperref}
\usepackage{enumitem}
\usepackage{listings}
\usepackage{titlesec}
\usepackage{fancyhdr}
\usepackage[protrusion=true,expansion=false]{microtype}

\definecolor{Ink}{HTML}{1B2430}
\definecolor{Muted}{HTML}{5B6472}
\definecolor{Accent}{HTML}{0077B6}
\definecolor{AccentTwo}{HTML}{D95F02}
\definecolor{Soft}{HTML}{F5F7FA}
\definecolor{SoftBlue}{HTML}{EAF5FB}
\definecolor{SoftOrange}{HTML}{FFF2E8}

\renewcommand{\familydefault}{\sfdefault}
\setlength{\parindent}{0pt}
\setlength{\parskip}{6pt}
\setlength{\emergencystretch}{4em}
\setlist{itemsep=2pt,topsep=4pt}

\hypersetup{
  colorlinks=true,
  linkcolor=Accent,
  urlcolor=Accent,
  citecolor=Accent,
  pdftitle={QuadrupolarDFT User Manual},
  pdfauthor={MRSpinDynamics project},
  pdfsubject={Beginner-friendly EFG and NQR calculation workflows with ABINIT and Elk}
}

\lstdefinestyle{shellstyle}{
  basicstyle=\ttfamily\small,
  backgroundcolor=\color{Soft},
  columns=fullflexible,
  keepspaces=true,
  breaklines=true,
  frame=leftline,
  framerule=2pt,
  rulecolor=\color{Accent},
  xleftmargin=0.8em,
  framesep=0.6em,
  showstringspaces=false
}
\lstset{style=shellstyle}
\newcommand{\code}[1]{{\small\nolinkurl{#1}}}
\newcommand{\cmd}[1]{\texttt{\small #1}}
\newcommand{\guidebox}[2]{%
  \begin{center}
  \setlength{\fboxsep}{9pt}%
  \fcolorbox{#1}{Soft}{\begin{minipage}{0.91\linewidth}#2\end{minipage}}
  \end{center}}
\newcommand{\chapterpath}[1]{%
  \guidebox{Accent}{\textbf{In this chapter.} #1}}

\titleformat{\chapter}[block]
  {\huge\bfseries\color{Ink}}{\color{Accent}\thechapter}{0.75em}{}
\titlespacing*{\chapter}{0pt}{-0.15in}{1.1em}
\titleformat{\section}{\Large\bfseries\color{Ink}}{\thesection}{0.6em}{}
\titleformat{\subsection}{\large\bfseries\color{Ink}}{\thesubsection}{0.6em}{}
\titleformat{\part}[display]
  {\centering\Huge\bfseries\color{Ink}}{\color{Accent}\Large PART \thepart}{0.5em}{}

\pagestyle{fancy}
\fancyhf{}
\lhead{\textcolor{Muted}{QuadrupolarDFT}}
\rhead{\textcolor{Muted}{User Manual}}
\cfoot{\textcolor{Muted}{\thepage}}
\renewcommand{\headrulewidth}{0.3pt}

\newcommand{\Cq}{C_{\mathrm{Q}}}
\newcommand{\Vzz}{V_{zz}}

\begin{document}
\pagenumbering{gobble}
\begin{titlepage}
\pagecolor{Soft}
\color{Ink}
\vspace*{1.1in}
{\Huge\bfseries QuadrupolarDFT\par}
\vspace{0.25in}
{\Large User Manual\par}
\vspace{0.12in}
{\large A practical route from crystal structure to NQR parameters\par}
\vspace{0.45in}
{\color{Accent}\rule{0.82\linewidth}{3pt}\par}
\vspace{0.35in}
\begin{minipage}{0.80\linewidth}
\large
This manual leads you from a crystal structure to electric-field-gradient (EFG)
tensors, quadrupolar coupling constants, and zero-field NQR transition
frequencies. It begins with a first static calculation, then builds toward
finite-temperature, displacement, and strain workflows. ABINIT provides the
efficient production path; Elk provides an all-electron accuracy check when the
near-nucleus EFG needs closer scrutiny.
\end{minipage}
\vfill
{\large MRSpinDynamics project \textbullet{} July 2026\par}
\end{titlepage}
\nopagecolor
\clearpage
\pagenumbering{roman}
\tableofcontents
\clearpage
\pagenumbering{arabic}

\part{Getting Started}

% =====================================================================
\chapter{Start Here: From a Structure to One NQR Prediction}
\label{ch:start}

\chapterpath{You will learn what the workspace does, which result to calculate
first, and how to check that a result is ready for scientific use. No prior DFT
workflow experience is assumed.}

The shortest useful QuadrupolarDFT workflow has three steps:
\begin{enumerate}
  \item start from a crystal structure whose origin and temperature are known;
  \item calculate an EFG tensor with ABINIT; and
  \item convert that tensor into $\Cq$, $\eta$, and NQR transition frequencies.
\end{enumerate}
That first static prediction is the foundation for every later calculation in
this manual. Do it before attempting phonons, temperature corrections, or strain
derivatives.

\section{This manual and the Beginner's Guide}

This is a \emph{task-oriented} manual. It tells you which files to prepare,
which commands to run, what the output means, and where to stop and check your
work. The companion
\code{docs/DFT\_for\_Spin\_Dynamics\_Beginners\_Guide.pdf} explains the ideas
behind those tasks: self-consistent fields, $k$-point meshes, basis sets,
pseudopotentials, all-electron methods, phonons, and parallel execution.

You do not need to read the whole Beginner's Guide first. Follow this manual and
use the Guide when a concept is unfamiliar. In particular:
\begin{itemize}
  \item read its introductory DFT and SCF sections before editing an ABINIT
    input for the first time;
  \item read its convergence discussion before trusting a reported EFG;
  \item read its PAW-versus-all-electron discussion before comparing ABINIT and
    Elk; and
  \item read its phonon and ABINIT DDB/ANADDB sections before Chapter
    \ref{ch:dfpt}.
\end{itemize}

\section{The four objects you will work with}

Keeping four objects separate makes the workflow much easier to reason about:
\begin{description}[style=nextline]
  \item[Structure]
  Lattice vectors, atomic coordinates, species, and the experimental or
  relaxed geometry from which they came.
  \item[Electronic-structure calculation]
  The ABINIT or Elk input, pseudopotentials or all-electron basis controls,
  exchange--correlation functional, $k$-point mesh, and convergence thresholds.
  \item[EFG tensor]
  A symmetric $3\times3$ tensor at a particular nucleus. This is the direct DFT
  result and should be retained, not only its principal value.
  \item[Magnetic-resonance observables]
  $\Cq$, $\eta$, principal axes, and transition frequencies derived using an
  isotope's spin and quadrupole moment.
\end{description}
If a prediction is surprising, this separation tells you where to look:
structure, electronic structure, tensor interpretation, or nuclear data.

\section{Install and run a tensor-only first example}

From the \code{QuadrupolarDFT} directory, create and activate a virtual
environment, then install the workspace:
\begin{lstlisting}
python -m venv .venv
# Linux/macOS:
. .venv/bin/activate
# Windows PowerShell:
# .venv\Scripts\Activate.ps1
python -m pip install -e .
\end{lstlisting}
The package can analyze a tensor without ABINIT or Elk. Save the following as
\code{quick\_start.py}; it is a useful installation check:
\begin{lstlisting}
from quadrupolar_dft import (
    EFGTensor, coupling_constant_hz, nqr_frequencies_hz
)

efg = EFGTensor.from_components([
    [-0.236543, 0.864057, 0.0],
    [0.864057, -0.236543, 0.0],
    [0.0, 0.0, 0.473086],
])
cq = coupling_constant_hz(
    efg.vzz_si, quadrupole_moment_barns=-0.02558
)
print("eta =", efg.eta)
print("Cq (MHz) =", cq / 1e6)
print("lines (Hz) =", nqr_frequencies_hz(
    spin=2.5, cq_hz=cq, eta=efg.eta
))
\end{lstlisting}
Run it with \cmd{python quick\_start.py}. The tensor is interpreted in ABINIT
atomic units by default; this convention is explained in Chapter
\ref{ch:install}.
For the complete worked calculation, use the provided NaNO$_2$ input in
\code{examples/abinit/}; Chapter \ref{ch:nano2} walks through the files and the
expected checks.

\guidebox{AccentTwo}{\textbf{What this command does not prove.} A successful
Python example checks parsing and tensor-to-NQR conversion. It does not show
that a DFT basis, $k$-point mesh, structure, or pseudopotential is accurate.
Those are physical convergence questions, not software installation tests.}

\section{Choose the smallest workflow that answers your question}

\begin{center}
\begin{tabular}{>{\raggedright\arraybackslash}p{0.25\linewidth}
                >{\raggedright\arraybackslash}p{0.30\linewidth}
                >{\raggedright\arraybackslash}p{0.34\linewidth}}
\toprule
Question & Start with & Add only when needed \\
\midrule
What lines should this structure produce? & Static ABINIT EFG & Elk cross-check
or finite-temperature correction \\
Why does calculation disagree with experiment? & Structure audit and ABINIT
convergence & Elk on the same geometry and functional \\
How do lines change with temperature? & Converged static reference & ABINIT
DFPT and finite displacements \\
How does strain drive or shift a transition? & Converged static reference &
Thirteen-job finite-strain set \\
\bottomrule
\end{tabular}
\end{center}

\section{A result is ready only after four checks}

Before quoting a frequency, confirm:
\begin{enumerate}
  \item \textbf{Identity:} the reported tensor belongs to the intended atom and
  isotope.
  \item \textbf{Numerics:} total energy and EFG components are stable to tighter
  cutoffs, denser $k$-point sampling, and stricter SCF convergence.
  \item \textbf{Conventions:} principal values are ordered by magnitude and the
  sign and units of $Q$, $\Vzz$, and $\Cq$ are recorded.
  \item \textbf{Provenance:} structure, code version, functional,
  pseudopotential or all-electron settings, and input/output files are saved.
\end{enumerate}
The remaining chapters turn those checks into repeatable workflows.

% =====================================================================
\chapter{Choosing a Calculation Strategy: Accuracy, Time, and Evidence}
\label{ch:strategy}

\chapterpath{You will decide when fast PAW calculations are sufficient, when an
all-electron check is worth its cost, and how to interpret disagreement between
ABINIT, Elk, and experiment.}

EFGs are unusually sensitive to the electron density close to a nucleus. That
makes the choice between a PAW calculation and an all-electron calculation more
important here than it may be for total energies or relaxed lattice constants.
In this repository, ABINIT and Elk therefore have complementary roles.

\section{ABINIT and Elk are different tools, not rival answers}

\begin{center}
\small
\begin{tabular}{>{\raggedright\arraybackslash}p{0.16\linewidth}
                >{\raggedright\arraybackslash}p{0.35\linewidth}
                >{\raggedright\arraybackslash}p{0.35\linewidth}}
\toprule
 & ABINIT PAW & Elk LAPW all-electron \\
\midrule
Best role & Production calculations, convergence surveys, relaxation, DFPT,
temperature ensembles, and strain derivatives & High-fidelity static
reference for selected structures and diagnosis of PAW/core sensitivity \\
Why it is faster & Smooth pseudo-wavefunctions and efficient $k$-point
parallelism & --- \\
Main accuracy risk & The PAW dataset may reconstruct the near-nucleus density
poorly for the EFG & Basis, muffin-tin radii, local orbitals, and $k$ points
must still be converged \\
Main time cost & Many moderate-cost solves in advanced workflows & Dense
generalized eigenproblems whose cost rises steeply with the LAPW basis \\
Repository support & Complete static, relaxation, phonon, displacement, and
strain path & Static reference and convergence examples \\
\bottomrule
\end{tabular}
\end{center}

The label \emph{all-electron} does not mean \emph{automatically exact}. Elk
still depends on the crystal structure, functional, $k$-point mesh, LAPW
cutoff, muffin-tin radii, local orbitals, and SCF tolerance. Its value is that it
removes the PAW dataset as one major source of uncertainty.

\section{Budget the workflow in solves, not a universal wall time}

A statement such as \emph{Elk is ten times slower} would be misleading: timing
depends strongly on structure, hardware, linear-algebra libraries, parallel
layout, and convergence settings. A more portable estimate counts independent
electronic-structure solves:
\begin{itemize}
  \item a static prediction needs roughly one converged solve, plus cutoff and
    $k$-point checks;
  \item a convergence survey needs a baseline and every off-baseline point;
  \item a finite-displacement temperature calculation needs the phonon work,
    one equilibrium EFG, and two EFG calculations for each retained mode;
  \item the six-component central finite-strain workflow needs thirteen EFG
    calculations: one equilibrium and positive/negative strain for each
    component; and
  \item an Elk check should normally cover a few representative structures,
    with its own \code{rgkmax} and $k$-point checks.
\end{itemize}
This arithmetic explains the repository strategy: use ABINIT for broad or
derivative-heavy work, and spend the more expensive all-electron calculation
where it changes the strength of the scientific conclusion.

\section{A practical two-tier protocol}

\begin{enumerate}
  \item \textbf{Audit the structure.} Record its source, temperature,
  occupancy, symmetry choices, and whether it is experimental or relaxed.
  \item \textbf{Converge the ABINIT EFG.} Vary cutoff and $k$-point density,
  monitoring the full tensor and $\eta$, not only total energy.
  \item \textbf{Compare with the available evidence.} Consider experimental
  uncertainty, temperature, polymorph, isotope, and line assignment.
  \item \textbf{Run Elk on the same geometry and functional} when the result is
  publication-critical, unexpectedly discrepant, or sensitive to a small
  difference between transverse EFG components.
  \item \textbf{Escalate selectively.} Use the diagnosis below to decide
  whether to improve the PAW dataset, structure, functional, or thermal model.
\end{enumerate}

\section{What the NaNO\texorpdfstring{$_2$}{2} comparison teaches}

The repository's NaNO$_2$ example holds the PBE functional and ICSD 82857
geometry fixed. The measured and calculated nitrogen results are:
\begin{center}
\begin{tabular}{lcc}
\toprule
Source & $\eta$ & $|\Cq|$ (MHz) \\
\midrule
Experiment & approximately 0.38 & 5.49 \\
ABINIT PAW & 0.112 & 5.03 \\
Elk LAPW, all-electron & approximately 0.333 & approximately 5.94 \\
\bottomrule
\end{tabular}
\end{center}
Elk's much better asymmetry on the same structure and functional points to the
selected nitrogen PAW dataset's reconstruction of the transverse,
near-nucleus EFG as the main limitation. It does \emph{not} prove that Elk will
always be closer to experiment, nor that thermal and structural effects are
unimportant in other systems.

The Elk convergence example is also a time warning. Results were already
stable between \code{rgkmax=7} and 8 and between the tested $k$ meshes, while a
larger automatic basis became both much more expensive and numerically
problematic after muffin-tin radii changed. Use optimized BLAS and no more MPI
ranks than useful irreducible $k$ points; consult the Beginner's Guide for why
that parallel layout matters.

\section{Diagnose disagreement before doing more computation}

\begin{description}[style=nextline]
  \item[ABINIT and Elk disagree on the same geometry and functional.]
  Suspect PAW reconstruction or incomplete convergence first. Converge both
  representations and, if possible, test an alternative PAW dataset.
  \item[ABINIT and Elk agree, but both disagree with experiment.]
  Look next at structure, polymorph, temperature, disorder, exchange--correlation
  functional, nuclear quadrupole moment, and spectral assignment.
  \item[$|\Cq|$ is stable but $\eta$ is not.]
  Inspect the tensor components directly. When $V_{xx}$ and $V_{yy}$ are close,
  a small absolute error can cause a large relative change in $\eta$.
  \item[Only the sign disagrees.]
  NQR transition frequencies often do not determine the sign of $\Cq$. Check
  tensor and nuclear-moment conventions before treating it as a physical
  discrepancy.
\end{description}

\guidebox{Accent}{\textbf{Decision rule.} Use ABINIT by default. Add Elk when a
near-nucleus accuracy question could change the interpretation, and compare on
the same structure and functional. The useful result is not which code wins,
but which source of uncertainty the comparison isolates.}

% =====================================================================
\chapter{Scope and Physical Background}
\label{ch:scope}

\chapterpath{You will connect the calculated EFG tensor to quadrupolar
Hamiltonians and see why a static 0\,K result may need a temperature
correction.}

QuadrupolarDFT addresses one link in the magnetic-resonance modelling chain: it
turns an \emph{ab initio} electronic-structure result into the parameters that
control a nuclear quadrupole resonance (NQR) experiment.

Readers new to electronic-structure theory may want to start with the companion
\code{docs/DFT\_for\_Spin\_Dynamics\_Beginners\_Guide.pdf}, which builds up the
background this manual assumes --- the SCF loop, $k$-points and basis sets, PAW
versus all-electron treatments, and phonons --- from a spin-dynamics standpoint.

A nucleus with spin $I \geq 1$ carries an electric quadrupole moment $Q$ that
couples to the electric-field-gradient (EFG) tensor $V_{ij} = \partial^2
V/\partial x_i \partial x_j$ produced by the surrounding charge distribution. The
EFG is symmetric and, away from the nucleus, traceless (Laplace's equation), so
in its principal-axis system it is described by two numbers: the largest-magnitude
component $\Vzz$ and the asymmetry parameter
\begin{equation}
\eta = \frac{V_{xx} - V_{yy}}{\Vzz}, \qquad |\Vzz| \geq |V_{yy}| \geq |V_{xx}|,
\qquad 0 \leq \eta \leq 1 .
\end{equation}
The quadrupolar coupling constant follows from $\Vzz$ and the nuclear quadrupole
moment,
\begin{equation}
\Cq = \frac{e\,Q\,\Vzz}{h},
\label{eq:cq}
\end{equation}
and the zero-field NQR transition frequencies follow from $\Cq$, $\eta$ and the
spin $I$ by diagonalizing the quadrupole Hamiltonian
\begin{equation}
\frac{H_{\mathrm Q}}{h} = \frac{\Cq}{4 I (2I-1)}
  \left[\,3 I_z^2 - I(I+1) + \eta\,(I_x^2 - I_y^2)\,\right].
\label{eq:hq}
\end{equation}
For spin $I=1$ (e.g.\ $^{14}$N) this gives the familiar triplet
\begin{equation}
\nu_{\pm} = \tfrac{3}{4}\Cq\left(1 \pm \tfrac{\eta}{3}\right), \qquad
\nu_{0} = \nu_{+} - \nu_{-} = \tfrac{1}{2}\Cq\,\eta .
\label{eq:spin1}
\end{equation}

\section{Why a finite-temperature correction is needed}

A static DFT calculation evaluates $V_{ij}$ at a single, fixed nuclear
configuration at $0$\,K. Measured NQR frequencies, however, are taken at finite
temperature and are strongly temperature dependent: in molecular crystals the
lines typically \emph{fall} as temperature rises (the Bayer law). Two physically
distinct effects separate the static DFT number from the measurement:

\begin{enumerate}
  \item \textbf{Thermal expansion / geometry.} The lattice constants and internal
    coordinates at temperature $T$ differ from the relaxed $0$\,K cell. This is a
    static effect: it is captured by evaluating the EFG on a
    temperature-appropriate structure.
  \item \textbf{Vibrational averaging.} The nucleus samples a thermal
    distribution of positions through the phonons. The measured EFG is the tensor
    averaged over that motion, including zero-point motion, which is present even
    at $0$\,K. For molecular crystals this dynamic term usually dominates:
    low-frequency librations reorient the EFG principal axes and reduce the line
    frequencies.
\end{enumerate}

QuadrupolarDFT provides the static layer (Chapter~\ref{ch:static}) and a harmonic
treatment of the vibrational term (Chapters~\ref{ch:finiteT}--\ref{ch:workflow}).

% =====================================================================
\chapter{Installation and Conventions}
\label{ch:install}

\chapterpath{You will install the Python helpers and fix the tensor, unit, and
nuclear-data conventions that every later result depends on.}

\section{Installation}

The package depends only on NumPy. From the \code{QuadrupolarDFT/} directory:
\begin{lstlisting}
python -m pip install -e .
python -m pip install -e ".[dev]"     # adds ruff for linting
python -m unittest discover -s tests
\end{lstlisting}
On a system whose Python is externally managed (PEP~668), the package can also be
used without installation by setting \code{PYTHONPATH} to the \code{src/}
directory, since NumPy is the only requirement:
\begin{lstlisting}
export PYTHONPATH="$PWD/src"
python3 examples/abinit/efg_temperature.py --help
\end{lstlisting}

\section{Tensor and unit conventions}

\code{EFGTensor} stores the symmetric tensor in SI units (V/m$^2$) and sorts the
principal components as $|\Vzz| \geq |V_{yy}| \geq |V_{xx}|$. It accepts input in
ABINIT atomic units (the default), in SI, or in the $10^{21}$\,V/m$^2$ scale, and
exposes \code{vzz\_si} and \code{eta}. Tensor averaging
(\code{average\_tensors}) operates on the tensor components and only then
diagonalizes the result, so that asymmetry information survives the average.

\section{Link to the spin-dynamics simulator}

The companion \code{spin\_dynamics} package parameterizes a quadrupolar site by
its $\eta=0$ transition frequency $\nu_{\mathrm Q}$. This is related to $\Cq$ by
the prefactor of Eq.~\eqref{eq:hq}:
\begin{equation}
\nu_{\mathrm Q} = \Cq \cdot \frac{d}{4 I (2I-1)} =
\begin{cases}
\tfrac{3}{4}\,\Cq & I = 1, \\[4pt]
\tfrac{1}{2}\,\Cq & I = \tfrac{3}{2},
\end{cases}
\end{equation}
with $d = 3$ for spin~1 and $d = 6$ for spin~3/2. The cross-project
\code{mr\_integration} layer uses this mapping to feed DFT parameters into the
simulator and to compare predictions against the measured NQR database.

\part{Core Calculations}

% =====================================================================
\chapter{Static EFG Workflow}
\label{ch:static}

\chapterpath{You will prepare a production ABINIT EFG calculation, parse its
tensor, convert it to NQR parameters, and test numerical convergence.}

The static path produces $\Cq$, $\eta$ and the NQR lines from an ABINIT EFG
calculation.

\begin{enumerate}
  \item \code{format\_abinit\_efg\_block} emits the \code{nucefg}/\code{quadmom}
    input lines that request an EFG calculation (PAW datasets are required;
    norm-conserving pseudopotentials are not suitable).
  \item ABINIT is run on the input (see \code{examples/abinit/run\_nano2\_efg\_wsl.sh}).
  \item \code{parse\_abinit\_efg} parses the \code{Electric Field Gradient
    Calculation} blocks of the output into \code{AbinitEFGRecord} objects, one per
    atom, carrying $\Cq$, $\eta$, the principal values/axes, and the total EFG
    tensor.
  \item \code{coupling\_constant\_hz} and \code{nqr\_frequencies\_hz} convert a
    tensor and a nuclear quadrupole moment into $\Cq$ and the unique zero-field
    transition frequencies (Eqs.~\eqref{eq:cq}--\eqref{eq:hq}).
\end{enumerate}

\code{examples/parse\_abinit\_efg.py} demonstrates parsing an output file. The EFG
components, not the total energy, must be converged against \code{ecut},
\code{pawecutdg} and the $k$-point mesh.

\guidebox{Accent}{\textbf{Static-workflow checkpoint.} Save the complete tensor,
atom index, isotope and quadrupole moment, ABINIT input/output, PAW dataset
identifiers, code version, structure source, and a convergence table. A single
printed $\Cq$ is not a reproducible DFT result.}

% =====================================================================
\chapter{Finite-Temperature EFG}
\label{ch:finiteT}

\chapterpath{You will see how phonon motion changes a static EFG and how the
package converts mode curvatures into temperature-dependent NQR lines.}

This chapter explains the model and its outputs. If you only need to understand
a supplied set of mode curvatures, read it straight through. If you need to
calculate those curvatures, continue to the staged workflow in Chapter
\ref{ch:dfpt}. The Beginner's Guide gives the more detailed phonon background.

\section{Harmonic vibrational averaging}

Expanding the EFG tensor to second order in the mass-weighted phonon normal
coordinates $Q_k$ and averaging over the harmonic distribution gives
\begin{equation}
\langle V_{ij}\rangle(T) = V_{ij}^{\mathrm{eq}}
  + \tfrac{1}{2}\sum_k \frac{\partial^2 V_{ij}}{\partial Q_k^2}\,
  \langle Q_k^2\rangle(T),
\label{eq:avg}
\end{equation}
where the mean-square amplitude of mode $k$ (angular frequency $\omega_k$) is
\begin{equation}
\langle Q_k^2\rangle(T) = \frac{\hbar}{2\omega_k}
  \coth\!\left(\frac{\hbar\omega_k}{2 k_{\mathrm B} T}\right).
\label{eq:q2}
\end{equation}
The $\coth$ factor equals $2 n(\omega_k,T) + 1$ with $n$ the Bose--Einstein
occupation. It tends to $1$ as $T\to 0$ (pure zero-point motion --- which is why
the static EFG is never exactly what is measured) and to
$2 k_{\mathrm B} T / \hbar\omega_k$ in the classical high-temperature limit. These
quantities live in \code{quadrupolar\_dft.thermal}.

\textbf{Tensor-frame averaging.} The average in Eq.~\eqref{eq:avg} is taken on the
tensor components in a fixed crystal frame; the result is diagonalized only
afterwards. Librational motion reorients the principal-axis system, so averaging
$\Vzz$ or $\eta$ directly would discard exactly that effect. As a consequence
$\eta$ is itself temperature dependent. \code{thermally\_averaged\_efg} performs
the average and removes any residual trace (the EFG is traceless by Laplace's
equation; the finite-difference curvatures of Section~\ref{sec:fd} amplify
rounding noise into a spurious trace, which is projected out).
\code{efg\_temperature\_sweep} returns $\Cq(T)$, $\eta(T)$ and the lines at each
requested temperature.

\section{The analytic Bayer limit}

Where per-mode EFG curvatures are not available, the temperature \emph{form} can
be validated directly against a measured line series with the single-libration
Bayer--Kushida model
\begin{equation}
\nu(T) = \nu_0\left[1 - a\,\coth\!\left(\frac{\hbar\omega}{2 k_{\mathrm B}T}\right)\right],
\label{eq:bayer}
\end{equation}
where $a$ is a small dimensionless librational amplitude.
\code{fit\_bayer\_single\_mode} fits $(\nu_0, a, \omega)$ to measured $\nu(T)$
data with a NumPy-only routine: for a fixed $\omega$ the model is linear in
$\nu_0$ and $\nu_0 a$, so a grid scan over $\omega$ with a linear least-squares
solve at each point suffices.

\guidebox{AccentTwo}{\textbf{Interpretation limit.} A good Bayer fit shows that a
line series has the expected thermal form. It does not identify a unique phonon
mode or replace a converged first-principles vibrational calculation.}

\part{Advanced Workflows}

% =====================================================================
\chapter{The DFPT Finite-Displacement Workflow}
\label{ch:workflow}
\label{ch:dfpt}

\chapterpath{You will run the staged ABINIT phonon and displacement workflow,
collect its EFG tensors, and validate the numerical derivatives before using
them for temperature prediction.}

The mode curvatures $\partial^2 V_{ij}/\partial Q_k^2$ in Eq.~\eqref{eq:avg} come
from displacing the structure along each phonon mode and recomputing the EFG.
Because the DFT runs are expensive and run outside the package, the workflow
separates into three stages with a calculation between each, preceded by an
optional structure-relaxation stage (Stage~0) that puts the geometry at an energy
minimum first. The CLI \code{examples/abinit/efg\_temperature.py} drives all of
them.

\begin{center}
\begin{tabular}{cll}
\toprule
Stage & Python action & Expensive external calculation \\
\midrule
0 (optional) & write/collect relaxation & ABINIT relaxation \\
1 & write phonon inputs & ABINIT DFPT, then ANADDB \\
2 & write displaced inputs & one ABINIT EFG per job \\
3 & collect and average & none \\
\bottomrule
\end{tabular}
\end{center}

\guidebox{Accent}{\textbf{Before Stage 1.} Finish the static convergence study
first. Phonons and displacements multiply the cost of every weak choice in the
base input. Copy the converged input into a new run directory and preserve it
unchanged as the workflow reference.}

\section{Normal-coordinate convention}
\label{sec:fd}

A phonon eigenvector $\boldsymbol\varepsilon_k$ is mass-weighted and normalized
($\sum |\varepsilon_k|^2 = 1$). A normal-coordinate step $\delta Q$ displaces atom
$i$ (mass $m_i$) by
\begin{equation}
\Delta \mathbf r_i = \delta Q\, \frac{\boldsymbol\varepsilon_{k,i}}{\sqrt{m_i}},
\end{equation}
and the curvature is the central difference of the EFG tensors at $0, \pm\delta Q$,
\begin{equation}
\frac{\partial^2 V_{ij}}{\partial Q_k^2} \approx
  \frac{V_{ij}(+\delta Q) - 2 V_{ij}(0) + V_{ij}(-\delta Q)}{\delta Q^2}.
\end{equation}
These units are consistent with $\langle Q_k^2\rangle$ from Eq.~\eqref{eq:q2}, so
the product in Eq.~\eqref{eq:avg} is an EFG.
\code{normal\_coordinate\_step\_si} chooses $\delta Q$ from a target maximum
Cartesian displacement (default $0.05$\,\AA).

\section{Stage 0 --- structure relaxation (optional)}
\label{sec:relax}

The harmonic phonon model is meaningful around a mechanically stable reference.
If the chosen structure has substantial residual forces, imaginary
$\Gamma$-point modes and biased equilibrium parameters can follow. An
experimental structure may still be the right static comparison, so retain it
and document any relaxed structure as a distinct model rather than silently
replacing one with the other. \cmd{efg\_temperature.py relax} strips the EFG
keywords from a converged static input and prepends a BFGS ionic-relaxation block
(\cmd{ionmov 2}, \cmd{optcell 0}, \code{tolmxf}); \code{run\_relax\_wsl.sh} runs
ABINIT. \cmd{efg\_temperature.py relax-collect} then reads the relaxed geometry
from the output footer (\code{-outvars: echo values of variables after
computation}, the same block abipy reads) and writes a relaxed static EFG input.

\begin{lstlisting}
python3 examples/abinit/efg_temperature.py relax \
    --base examples/abinit/nano2_efg.abi --out runs/nano2_relax
bash examples/abinit/run_relax_wsl.sh runs/nano2_relax
python3 examples/abinit/efg_temperature.py relax-collect \
    --base examples/abinit/nano2_efg.abi --abo runs/nano2_relax/relax.abo \
    --out runs/nano2_relax
\end{lstlisting}

The product, \code{relaxed.abi}, is an ordinary static EFG input (EFG keywords
intact) at the relaxed geometry; pass it as \code{--base} to Stages~1--2 in place
of the hand-written input. The relaxation is internal-coordinates-only: the cell
is held at the input value (\cmd{optcell 0}). For molecular crystals this is
usually preferable to a full cell relaxation, which under GGA tends to over-expand
the lattice and can worsen the EFG relative to the measured geometry.
Temperature-dependent cell expansion is not included in the present harmonic
finite-displacement workflow; treat it as a separate structural contribution.

\section{Stage 1 --- phonons}

\cmd{efg\_temperature.py phonon} writes a two-dataset ABINIT input (ground state
plus a $\Gamma$-point DFPT response) and the matching anaddb input, derived from a
converged static EFG input. \code{run\_phonon\_wsl.sh} runs ABINIT, selects the
DFPT derivative database (the largest \code{*\_DDB}), and runs anaddb to produce
the phonon frequencies and eigenvectors.

\begin{lstlisting}
python3 examples/abinit/efg_temperature.py phonon \
    --base examples/abinit/nano2_efg.abi --out runs/nano2_ph
bash examples/abinit/run_phonon_wsl.sh runs/nano2_ph
\end{lstlisting}

The generated DFPT input is a starting template. Two settings are required for
PAW~+~GGA phonons and are emitted automatically: \cmd{iscf2 7} (the response
driver requires \code{iscf}~$\leq 9$) and \cmd{pawxcdev 0} (the GGA exchange
correlation kernel is not implemented in the development scheme). Converge the
tolerances and $k$-mesh for production use.

\section{Stage 2 --- displaced EFG runs}

\cmd{efg\_temperature.py displace} parses the anaddb eigenvectors
(\code{parse\_anaddb\_modes}; acoustic and imaginary modes are dropped, and modes
are paired with frequencies by mode number so that gaps in the eigenvector output
do not misalign them), generates one displaced EFG input per $\pm$ mode plus a
manifest, and writes positions as fractional \code{xred} coordinates.
\code{run\_finite\_displacement\_wsl.sh} runs ABINIT on every input.

\begin{lstlisting}
python3 examples/abinit/efg_temperature.py displace \
    --base examples/abinit/nano2_efg.abi --anaddb runs/nano2_ph/anaddb.out \
    --target 2 --max-modes 6 --out runs/nano2_disp
bash examples/abinit/run_finite_displacement_wsl.sh runs/nano2_disp
\end{lstlisting}

\code{--target} is the $0$-based index of the resonant nucleus; \code{--max-modes}
restricts the calculation to the lowest-frequency librations, which dominate the
temperature dependence. If the anaddb eigenvector layout differs from the parser's
expectation, \cmd{--modes modes.json} supplies frequencies and eigenvectors
directly.

\section{Stage 3 --- collect}

\cmd{efg\_temperature.py collect} parses the displaced EFG outputs, central
differences the mode curvatures, runs the temperature sweep, and reports
$\Cq(T)$, $\eta(T)$, the lines, and $\mathrm d\nu/\mathrm dT$.

\begin{lstlisting}
python3 examples/abinit/efg_temperature.py collect \
    --workdir runs/nano2_disp --temperatures 0,77,150,300 --quadmom 0.02044
\end{lstlisting}

\guidebox{AccentTwo}{\textbf{Do not skip the derivative check.} Repeat selected
modes with a smaller and larger displacement. The resulting curvature should
be stable: too small a step amplifies SCF noise, while too large a step samples
anharmonic response. Inspect plus/minus tensor changes, imaginary modes, and
the residual trace before interpreting a temperature slope.}

\section{Running ABINIT in parallel}
\label{sec:parallel}

The phonon and EFG runs dominate the wall-clock cost, and ABINIT (MPI-built)
parallelizes them efficiently over $k$-points. All the runner scripts source a
small shared helper, \code{examples/abinit/abinit\_parallel.sh}, and honor the
\code{ABINIT\_NP} environment variable:

\begin{itemize}
  \item \code{ABINIT\_NP} unset (or \code{1}) --- serial.
  \item \code{ABINIT\_NP=}$N$ --- launch \cmd{mpirun -np $N$ abinit}.
  \item \code{ABINIT\_NP=auto} --- choose $N$ automatically (see below) and fall
    back to serial if no MPI runtime is found.
\end{itemize}

With the default $k$-point parallelism (\cmd{paral\_kgb 0}), efficiency is good
when the number of processes divides the number of $k$-points in the IBZ. The
helper function \code{abinit\_optimal\_np} encodes this: it returns the largest
divisor of $n_{\mathrm{kpt}}$ that is $\le$ the core count. For the NaNO$_2$ cell
($n_{\mathrm{kpt}}=36$) on a 20-core machine that is $N=18$ (two $k$-points per
process), which scales near-ideally ($\sim 18\times$). \code{auto} detects
$n_{\mathrm{kpt}}$ from an existing \code{.abo} or a quick \cmd{abinit --dry-run},
and the core count from \code{nproc}. Override the launcher with
\code{ABINIT\_MPIRUN} (e.g.\ \code{mpiexec}, or \cmd{mpirun --oversubscribe}).
The driver \code{nano2\_relaxation\_study.sh} uses \code{auto} by default; pass
\code{--np}~$N$ to pin a value or \cmd{--np 1} to force serial.

Parallelization changes only how ABINIT is launched, not the inputs, so it leaves
all computed quantities (and any relaxed-vs-unrelaxed comparison) unchanged.

\section{Convergence of \texorpdfstring{$\eta$}{eta} and the asymmetry gap}
\label{sec:convergence}

The coupling $\Cq \propto \Vzz$ tracks the largest EFG eigenvalue and converges
quickly, but the asymmetry $\eta = (V_{xx}-V_{yy})/\Vzz$ is a normalized
\emph{difference} of the two smaller eigenvalues. It therefore converges much more
slowly with the PAW double-grid cutoff \code{pawecutdg} and the $k$-mesh
\code{ngkpt}, and a calculation can reproduce $\Cq$ while still reporting a
too-axial (low) $\eta$ simply because the transverse tensor components are not yet
converged. \code{examples/abinit/efg\_convergence.py} isolates that effect: it
sweeps \code{ecut}, \code{pawecutdg}, and \code{ngkpt} one knob at a time around
the baseline read from a base input, then tabulates $\eta$ (and $\Cq$) per knob
with successive deltas and a converged flag.

\begin{lstlisting}
PYTHONPATH=src python3 examples/abinit/efg_convergence.py generate \
    --base examples/abinit/nano2_efg.abi --target 2 \
    --ecut 25,30,35 --pawecutdg 50,60,80,100 --ngkpt 4,6,8 \
    --out runs/nano2_conv
bash examples/abinit/run_convergence_efg_wsl.sh runs/nano2_conv
PYTHONPATH=src python3 examples/abinit/efg_convergence.py collect \
    --workdir runs/nano2_conv --out runs/nano2_conv/convergence.md \
    --csv runs/nano2_conv/convergence.csv
\end{lstlisting}

Each knob's ladder shares one baseline job (only off-baseline values are staged as
extra runs), so the baseline is computed once and anchored across all three
sweeps. \code{run\_convergence\_efg\_wsl.sh} runs each job in a WSL-native scratch
directory (a \code{mktemp} dir under \code{\$TMPDIR}, overridable with
\code{ABINIT\_SCRATCH\_DIR}) and copies only \code{<name>.abo} back into the
workdir. This matters on OneDrive-synced \code{/mnt/c} checkouts, where ABINIT
scratch and \code{.abo} I/O is slow and OneDrive file locks can stall a run so
that no \code{.abo} ever appears.

\paragraph{NaNO$_2$ result.} For the $^{14}$N site the sweep is decisive: $\eta$
is flat to $\sim 10^{-4}$ across \code{pawecutdg} ($50\!\to\!100$~Ha) and
\code{ngkpt} ($4^3\!\to\!8^3$, $n_{\mathrm{kpt}}$ up to $260$), while \code{ecut}
moves it only slightly and \emph{away} from experiment (knee near $30$~Ha).
Convergence --- and, because the inputs run with \cmd{nsym 1}, symmetry pinning
--- is therefore ruled out as the cause of the low $\eta$. The geometry is the
dominant lever instead:

\begin{center}
\begin{tabular}{lcc}
\toprule
geometry & $\eta$ & $\Cq$ (MHz) \\
\midrule
hand-built starter & 0.043 & $-5.17$ \\
ICSD 82857 (experimental) & 0.112 & $-5.03$ \\
experiment (77--300~K NQR) & $\sim 0.38$ & $\pm 5.49$ \\
\bottomrule
\end{tabular}
\end{center}

The experimental cell raises $\eta$ by $\sim 2.6\times$ over the starter geometry
but still under-predicts the measured value. The cause of that residual was
pinned down with an all-electron cross-check.

\paragraph{All-electron cross-check (Elk): the gap is PAW-sensitive.} The
transverse EFG components that set $\eta$ are a near-nucleus quantity, so they are
sensitive to how the pseudopotential reconstructs the core region. Repeating the
calculation with the \emph{same} PBE functional and the \emph{same} ICSD geometry
in the all-electron full-potential LAPW code Elk --- the reference method for EFGs
--- gives a very different $\eta$:

\begin{center}
\begin{tabular}{lcc}
\toprule
method (PBE, ICSD geometry) & $\eta$ & $|\Cq|$ (MHz) \\
\midrule
experiment (77--300~K NQR) & $\sim 0.38$ & 5.49 \\
ABINIT --- PAW (Pseudodojo) & 0.112 & 5.03 \\
Elk --- all-electron LAPW & $\approx 0.333$ & 5.9 \\
\bottomrule
\end{tabular}
\end{center}

All-electron PBE recovers $\eta \approx 0.333$ (converged over \code{rgkmax}
$7\!\to\!8$ and the $k$-mesh, $18\!\to\!48$ points), next to the measured $0.38$,
whereas the PAW value sits $\sim 3\times$ lower. Same functional, same geometry:
the comparison isolates a strong \textbf{PAW/pseudopotential sensitivity}. The
selected Pseudodojo PBE nitrogen dataset under-represents the near-nucleus
transverse EFG in this case; geometry and functional do not explain the
ABINIT--Elk difference.
(Harmonic vibrational averaging, the finite-temperature workflow above, in fact
\emph{lowers} $\eta$, so dynamics does not close the gap upward.) The ABINIT-side
An ABINIT-side next test is therefore a harder or EFG-tuned nitrogen PAW dataset
before changing the functional. The Elk cross-check workflow --- input, an \code{EFG.OUT} parser,
and an MPI-parallel runner with a convergence watchdog --- is in
\code{examples/elk/}; its \code{README.md} records the inputs, convergence
results, performance observations, and interpretation.

% =====================================================================
\chapter{Finite-Strain EFG Drive Couplings}
\label{ch:strain}

\chapterpath{You will generate, collect, and validate the thirteen calculations
needed for central finite-strain EFG derivatives and transition-drive
couplings.}

Piezoelectric NQR drive needs a different DFT derivative from the
finite-temperature workflow.  Instead of displacing atoms along optical phonon
normal modes, the experiment applies a macroscopic strain field through the
piezoelectric crystal.  To first order,
\[
  V_{ij}(\boldsymbol\epsilon)
    \simeq V_{ij}(0)
      + \sum_{\mu} \frac{\partial V_{ij}}{\partial \epsilon_{\mu}}\,
        \epsilon_{\mu},
\]
where $\mu \in \{xx,yy,zz,yz,xz,xy\}$ is the symmetric strain component in Voigt
notation.  The transition-drive coefficient is obtained by projecting
$\partial V_{ij}/\partial\epsilon_\mu$ into the eigenbasis of the equilibrium
quadrupolar Hamiltonian.  This is the quantity used by the piezoelectric
drive/detection model in PythonSpinDynamics.

\section{Glycine polymorph guard}

The bundled glycine structure is CCDC~189379, space group \cmd{P 21}.  This
non-centrosymmetric polymorph is symmetry-allowed to be piezoelectric, whereas
centrosymmetric glycine polymorphs are not useful for an electric-field-driven
piezoelectric transducer.  The CLI therefore performs a CIF metadata check
before writing jobs:

\begin{lstlisting}
PYTHONPATH=src python3 examples/abinit/strain_efg_coupling.py check
\end{lstlisting}

For CCDC~189379, the expanded ABINIT cell contains two symmetry-related nitrogen
sites.  The zero-based target indices printed by \code{prepare-base} are
\code{1} and \code{11}; the example below uses index~\code{1}.

\section{Static and strained ABINIT jobs}

The strain workflow begins by preparing a static EFG input directly from the CIF.
On Ubuntu's packaged ABINIT~9.10.4 PAW set, the glycine H/C/N/O XML files are in
\code{Pseudodojo\_paw\_pw\_standard}.  Those PAW spheres overlap on short X--H
bonds, so the exploratory run used \cmd{pawovlp 25}; production calculations
should replace this with smaller-radius PAW datasets or a validated
pseudopotential choice.

\begin{lstlisting}
PYTHONPATH=src python3 examples/abinit/strain_efg_coupling.py prepare-base \
  --out runs/glycine_static/glycine_efg.abi \
  --pawovlp 25

bash examples/abinit/run_static_efg_wsl.sh \
  runs/glycine_static/glycine_efg.abi
\end{lstlisting}

After the static input is validated, the generator writes the equilibrium plus
central-difference strained structures:

\begin{lstlisting}
PYTHONPATH=src python3 examples/abinit/strain_efg_coupling.py generate \
  --base runs/glycine_static/glycine_efg.abi \
  --target-atom-index 1 \
  --out runs/glycine_strain

bash examples/abinit/run_strain_efg_wsl.sh runs/glycine_strain
\end{lstlisting}

The collection step parses each \code{.abo}, finite-differences the target
nucleus EFG tensor, diagonalizes the equilibrium spin-1 Hamiltonian, and writes
JSON/CSV coupling tables:

\begin{lstlisting}
PYTHONPATH=src python3 examples/abinit/strain_efg_coupling.py collect \
  --workdir runs/glycine_strain \
  --quadmom 0.02044 \
  --strain-peak 1e-5 \
  --json runs/glycine_strain/couplings.json \
  --csv runs/glycine_strain/couplings.csv
\end{lstlisting}

\section{First-pass glycine result}

The exploratory glycine calculation completed all 13 jobs (equilibrium plus
$\pm$ strains for six components).  The equilibrium EFG at the target nitrogen
was
\[
  (V_{xx},V_{yy},V_{zz})
    = (-1.9249,-1.2972,3.2221)\times 10^{21}\ \mathrm{V/m^2},
  \qquad \eta = 0.1948,
\]
corresponding to $\Cq = 1.5925$\,MHz for $^{14}$N.  The NQRDatabase glycine
amine-$^{14}$N entry is $\Cq=1.193$\,MHz and $\eta=0.528$, with observed lines at
737 and 1052\,kHz.  Thus the DFT run captures the MHz scale but makes the tensor
too axial; absolute frequencies are not yet production quality.

\begin{table}[htbp]
\centering
\begin{tabular}{llrrr}
\toprule
Transition & Best strain & Frequency (kHz) & Coupling (MHz/strain) &
Rabi at $10^{-5}$ strain (Hz) \\
\midrule
$1\to2$ & $xz$ & 155.124 & 1.456 & 7.28 \\
$0\to1$ & $yy$ & 1116.801 & 1.844 & 9.22 \\
$0\to2$ & $xz$ & 1271.924 & 2.189 & 10.95 \\
\bottomrule
\end{tabular}
\caption{Strongest projected quadrupolar drive coefficients from the first-pass
glycine finite-strain calculation.  The Rabi rate assumes a peak acoustic strain
of $10^{-5}$.}
\label{tab:glycine-strain-drive}
\end{table}

For the experimental two-line glycine model in PythonSpinDynamics, the strongest
DFT coefficients are mapped onto the observed lower and upper lines as practical
drive estimates: the lower 737\,kHz line uses the $yy$-driven $0\to1$ coefficient
($1.844$\,MHz/strain), and the upper 1052\,kHz line uses the $xz$-driven
$0\to2$ coefficient ($2.189$\,MHz/strain).  This keeps the detection model tied
to measured frequencies while using DFT-scale electric-field-gradient
derivatives.

\guidebox{AccentTwo}{\textbf{Reliability checkpoint.} These are exploratory
couplings, not design numbers. Before using them quantitatively, repeat the
workflow with converged \code{ecut}, \code{pawecutdg}, and \code{ngkpt}; relax
at least the hydrogen positions or use a neutron-quality structure; verify the
symmetry-related nitrogen index~\code{11}; test strain-step linearity; and
remove or quantify the effect of \cmd{pawovlp 25}.}

% =====================================================================
\chapter{Worked Example: \texorpdfstring{NaNO$_2$ $^{14}$N}{NaNO2 nitrogen-14}}
\label{ch:nano2}

\chapterpath{You will follow the repository's reference system from its source
structure through static and finite-temperature results, including the checks
that limit what may be concluded.}

Sodium nitrite is the reference case for the workspace: its $^{14}$N parameters
appear in the DFT runs, the spin-dynamics simulator, and the measured NQR
database. A six-mode finite-displacement calculation on the starter geometry
produces the temperature dependence in Table~\ref{tab:nano2}.

There are two useful ways to follow this example:
\begin{enumerate}
  \item use the small starter calculation to learn the mechanics of the staged
    workflow and recognize its failure checks; then
  \item use the ICSD 82857 structure, converged ABINIT calculation, and Elk
    comparison for the stronger accuracy argument.
\end{enumerate}
Do not combine the first track's convenient teaching numbers with the second
track's scientific conclusions.

\begin{table}[h]
\centering
\begin{tabular}{rrrl}
\toprule
$T$ (K) & $\Cq$ (MHz) & $\eta$ & lines (MHz) \\
\midrule
static & $-5.173$ & $0.043$ & 0.11, 3.82, 3.94 \\
0   & $-4.963$ & $0.016$ & 0.04, 3.70, 3.74 \\
77  & $-4.927$ & $0.014$ & 0.03, 3.68, 3.71 \\
150 & $-4.816$ & $0.002$ & 0.00, 3.61, 3.61 \\
300 & $-4.534$ & $0.033$ & 0.08, 3.36, 3.44 \\
\bottomrule
\end{tabular}
\caption{Finite-temperature $^{14}$N EFG of NaNO$_2$ from a six-mode harmonic
average on the starter geometry. The slope of the two strong lines is
$\mathrm d\nu/\mathrm dT \approx -1.0$ to $-1.1$\,kHz/K over $0$--$300$\,K.}
\label{tab:nano2}
\end{table}

The result reproduces the expected physics with no fitting: $|\Cq|$ and the lines
decrease smoothly with temperature (the Bayer sign), the $0$\,K value already lies
below the static one by about $4\%$ (zero-point motion), and $\eta$ shifts with
temperature. The slope, $\sim -1.0$\,kHz/K, has the correct sign and order of
magnitude relative to the measured coefficients of $-1.3$\,kHz/K ($\nu_-$) and
$-2.2$\,kHz/K ($\nu_+$).

\textbf{Why this is a teaching result, not a prediction.} The numbers above are
for the \emph{unrelaxed} starter geometry:
its static $\eta = 0.04$ disagrees with the experimental $0.38$, the two strong
lines come out near-degenerate ($\sim 3.7$\,MHz) rather than well split
($3.60$/$4.64$\,MHz), and the phonon spectrum contains imaginary modes. The
temperature trend is physically sensible, but the reference is not mechanically
stable. Run Stage~0 (\S\ref{sec:relax}) and pass \code{relaxed.abi} as the base
input to remove that particular defect. Relaxation alone does not make the
absolute EFG accurate; the later structure and Elk comparisons test the
remaining sources of error.

\textbf{Runnable comparison.} \code{examples/abinit/nano2\_relaxation\_study.sh}
automates exactly this check: it runs the full temperature workflow on both the
unrelaxed and the relaxed geometry and writes a side-by-side report (static EFG,
the sweep, and $\mathrm d\nu/\mathrm dT$) scored against the measured NaNO$_2$
$^{14}$N reference. See \code{examples/abinit/README\_relaxation\_study.md} for
the prerequisites and the one-command invocation.

\textbf{Relaxation outcome.} Running this on the starter geometry confirms the
pipeline behaves correctly and is qualitatively improved by relaxation, but also
exposes the starter model's limit. Relaxing the internal coordinates removes all imaginary
modes (the spectrum becomes dynamically stable, lowest libration $\sim 44$\,cm$^{-1}$)
and fixes the sign of every $\mathrm d\nu/\mathrm dT$, and it moves the asymmetry
and the line splitting toward experiment --- but only about a third of the way
($\eta: 0.04 \to 0.14$ against the measured $0.38$; strong-line splitting
$0.11 \to 0.37$\,MHz against $\sim 1.04$\,MHz). Because the cell was held fixed at
the hand-entered starter (\cmd{optcell 0}, unconverged \code{ecut}/\code{pawecutdg}/
\code{ngkpt}), this run cannot isolate the residual error. Quantitative work
therefore starts with a converged calculation
on an experimental structure (e.g.\ the ICSD~82857 input under
\code{structures/NaNO2/generated/}), followed by the ABINIT--Elk comparison in
Chapter~\ref{ch:strategy}. For the full slope magnitude, anharmonic (AIMD/PIMD)
averaging may also be needed near the ferroelectric transition.

\textbf{CIF validation outcome.} The bundled NaNO$_2$ CIFs now have a two-stage
check. \nolinkurl{examples/check_nano2_cifs.py} first ranks the available
structures by phase, temperature, occupancy and crystallographic metadata; for
the room-temperature ferroelectric NQR reference, ICSD~82857 is the best
structural target. \nolinkurl{examples/abinit/nano2_cif_efg_validation.sh} then stages the
top full-occupancy CIFs, runs static EFG calculations through the MPI-aware
ABINIT wrapper, and compares the resulting $^{14}$N lines against the measured
NaNO$_2$ reference. In the first validation sweep the best simulated line
agreement came from ICSD~15400, but the full-occupancy ferroelectric structures
all gave the same qualitative answer: $|\Cq|$ was close, while $\eta$ remained
too small ($\eta \simeq 0.11$--$0.12$ versus the measured $0.38$).

\textbf{Interpretation.} The NaNO$_2$ asymmetry error is not fixed by choosing a
different bundled CIF. The converged same-geometry, same-functional Elk
comparison then isolates the selected nitrogen PAW dataset as the dominant
source of the static ABINIT--Elk asymmetry gap. Thermal expansion, anharmonicity,
and structural disorder remain separate questions when comparing either static
calculation with finite-temperature experiment. High-temperature
\cmd{I m m m} CIFs should also be treated carefully: the available entries are
partial-occupancy/disordered models, and a literal static expansion can place
split sites nearly on top of each other, causing ABINIT PAW sphere-overlap
aborts. Use those only after constructing an ordered approximant, not as ordinary
static EFG inputs.

\guidebox{Accent}{\textbf{What the example establishes.} The software path
preserves tensor conventions and produces the expected thermal sign; relaxation
removes the starter structure's imaginary modes; structure screening does not
resolve the asymmetry; and the controlled Elk comparison reveals PAW
sensitivity. Each statement comes from a different check, and none should be
claimed from the temperature table alone.}

\part{Evidence and Reference}

% =====================================================================
\chapter{Validation Against the Measured Database}
\label{ch:validation}

\chapterpath{You will compare calculated static parameters and temperature
coefficients with measured records through the cross-project integration
layer.}

The cross-project \code{mr\_integration} layer closes the loop against the
NQRDatabase. \code{measured\_temperature\_coefficients} reads the stored
$\mathrm d\nu/\mathrm dT$ for a compound and isotope;
\code{compare\_temperature\_coefficients} matches a predicted set of
$(\text{frequency}, \mathrm d\nu/\mathrm dT)$ pairs --- from either a
\code{efg\_temperature\_sweep} (via \code{slopes\_from\_temperature\_points}) or a
Bayer fit --- to the measured coefficients by frequency, so that a DFT temperature
dependence can be checked against experiment. The same layer also validates the
static $(\Cq,\eta)$ against measured line positions and supplies the
$\Cq \leftrightarrow \nu_{\mathrm Q}$ conversion used by the simulator.

Use that comparison in the following order:
\begin{enumerate}
  \item select the compound, isotope, phase, and measurement temperature;
  \item compare static $\Cq$, $\eta$, and line ordering before comparing slopes;
  \item match calculated and measured lines by frequency, then inspect whether
    that assignment remains stable across temperature; and
  \item report residuals together with the DFT convergence and structural
    provenance, rather than treating the database value as an exact target.
\end{enumerate}
Agreement with experiment tests the complete model, not only the Python
conversion. Disagreement can originate in the measured assignment, structure,
nuclear data, electronic-structure approximation, or temperature model.

% =====================================================================
\chapter{Public API Reference}
\label{ch:api}

\chapterpath{Use this chapter after the narrative workflows when you need to
find the Python object or function that implements a particular step.}

\section{quadrupolar\_dft.tensors / quadrupolar / abinit}
\begin{itemize}
  \item \code{EFGTensor}, \code{average\_tensors} --- tensor conventions and
    frame-correct averaging.
  \item \code{coupling\_constant\_hz}, \code{nqr\_frequencies\_hz} --- $\Cq$ and
    zero-field transition frequencies.
  \item \code{AbinitEFGRecord}, \code{parse\_abinit\_efg},
    \code{format\_abinit\_efg\_block} --- ABINIT EFG input/output.
\end{itemize}

\section{quadrupolar\_dft.thermal}
\begin{itemize}
  \item \code{thermal\_quantum\_factor}, \code{bose\_occupation},
    \code{mean\_square\_normal\_coordinate}, \code{wavenumber\_to\_angular\_frequency}.
\end{itemize}

\section{quadrupolar\_dft.vibrational}
\begin{itemize}
  \item \code{VibrationalMode}, \code{efg\_curvature\_central\_difference},
    \code{thermally\_averaged\_efg}, \code{efg\_temperature\_sweep},
    \code{ThermalEFGPoint}.
  \item \code{bayer\_frequency}, \code{bayer\_slope\_hz\_per\_k},
    \code{fit\_bayer\_single\_mode}, \code{BayerFit}.
\end{itemize}

\section{quadrupolar\_dft.finite\_displacement / abinit\_phonon}
\begin{itemize}
  \item \code{Crystal}, \code{PhononMode}, \code{parse\_abinit\_structure},
    \code{generate\_displacement\_jobs}, \code{abinit\_input\_with\_positions},
    \code{write\_jobs}.
  \item \code{collect\_efg\_outputs}, \code{vibrational\_modes\_from\_collected},
    \code{vibrational\_modes\_from\_efg}, \code{efg\_tensor\_from\_record}.
  \item \code{phonon\_dfpt\_input}, \code{anaddb\_input},
    \code{parse\_anaddb\_modes}, \code{modes\_from\_arrays}.
\end{itemize}

\section{quadrupolar\_dft.strain\_coupling}
\begin{itemize}
  \item \code{generate\_strain\_jobs}, \code{write\_strain\_jobs} --- central
    finite-difference homogeneous strain inputs.
  \item \code{collect\_strain\_derivatives} --- parse strained EFG outputs and
    compute $\partial V_{ij}/\partial\epsilon_\mu$.
  \item \code{strain\_transition\_couplings} --- project strain derivatives onto
    spin-1 quadrupolar transition matrix elements.
  \item \code{glycine\_static\_efg\_input\_from\_cif},
    \code{cif\_structure\_metadata},
    \code{space\_group\_is\_likely\_centrosymmetric} --- glycine setup and
    polymorph guards for piezoelectric work.
\end{itemize}

\section{quadrupolar\_dft.relax}
\begin{itemize}
  \item \code{relax\_input} (ionic-relaxation input from a static EFG input),
    \code{parse\_relaxed\_structure} (relaxed geometry from the ABINIT output
    footer), \code{relaxed\_input} (relaxed static EFG input for Stages~1--2).
\end{itemize}

% =====================================================================
\chapter{Limits of the Current Workflow}
\label{ch:limits}

\chapterpath{You will distinguish implemented, validated capabilities from
known model limitations and choose an appropriate next validation step.}

The repository automates a substantial calculation chain, but it does not
automate scientific judgment. The following boundaries should appear in any
report that relies on the workflow:
\begin{itemize}
  \item The $\Cq \leftrightarrow \nu_{\mathrm Q}$ scale and several helpers are
    calibrated for spin~1 and spin~3/2; higher half-integer spins are not yet
    handled by the simulator-facing conversion.
  \item The generated DFPT input is a template; PAW DFPT settings, tolerances and
    $k$-mesh must be checked for production accuracy.
  \item The anaddb eigenvector parser targets ABINIT~9.x text output; a
    \code{--modes} JSON escape hatch covers other layouts.
  \item Finite-difference curvatures use the printed EFG precision; large
    displacement amplitudes or extra output digits may improve the second
    derivative, but step-size stability must still be demonstrated.
  \item The harmonic finite-temperature calculation does not include
    quasi-harmonic cell expansion or anharmonic sampling.
  \item Elk support is a static reference workflow, not a replacement backend
    for the ABINIT DFPT, displacement, and strain automation.
\end{itemize}

\section*{Choosing a method beyond these limits}

Use temperature-dependent experimental structures or a quasi-harmonic treatment
when cell expansion is important. Use AIMD/PIMD snapshot averaging when
anharmonicity or disorder invalidates the normal-mode picture. Convert other
phonon formats explicitly when the ANADDB text parser is unsuitable. In each
case, keep the downstream tensor conventions unchanged: average full tensors in
a common frame, then diagonalize and calculate the NQR observables.

\guidebox{Accent}{\textbf{A good stopping point.} A calculation is complete when
it answers the stated scientific question with documented numerical,
structural, and model uncertainty. Adding a more expensive method is useful
only when it tests an uncertainty that matters to that answer.}

\end{document}
